% !TeX root = ../main.tex
\section{Native code and intermediate representations}

If the goal is speed, why doesn't a compiler just translate the AST straight into
machine instructions in one pass? You can do that, and simple compilers do, but any
serious compiler inserts a stepping stone in between: an \textbf{intermediate
representation} (IR). This section is about why that extra layer exists and what it
buys you.

\subsection{why bother with an IR at all}

The AST is shaped like the \emph{source language}: it has \texttt{if} nodes, \texttt{while}
nodes, nested expressions, all the syntactic sugar a human programmer wanted to write.
Machine code is shaped like the \emph{target hardware}: registers, flags, jump
addresses, fixed-width instructions. Translating directly between two representations
that are both highly specific (one to the source language, one to the target chip)
means you'd need a separate translator for every (source language, target
architecture) pair, and you'd have to re-implement every optimization once per pair
too.

An IR sits in between: lower-level than the AST (closer to what a machine actually
does, one operation at a time, explicit temporaries instead of nested expressions) but
still machine-independent (no specific registers, no specific instruction encodings).
This buys you two things at once:
\begin{itemize}
  \item \textbf{Optimization} happens once, on the IR, and benefits every target
    architecture the compiler supports, instead of being duplicated per-architecture.
  \item \textbf{Retargeting} to a new architecture only requires writing a new
    \emph{backend} (IR $\to$ that architecture's machine code), reusing the entire
    front end (parsing, type checking, IR generation, optimization) unchanged.
\end{itemize}

\subsection{three-address code}

The canonical example IR is \textbf{three-address code} (TAC): every instruction has at
most one operator and at most three operands (typically: two operands and one
destination), and every intermediate result gets its own explicit name, a
\textbf{temporary}.

\begin{example}[Translating \texttt{x = a + b * c}]
The AST has a multiplication nested inside an addition. TAC flattens this into a
sequence of single-operation steps, introducing a temporary for the intermediate
product:
\begin{lstlisting}
t1 = b * c
x  = a + t1
\end{lstlisting}
Every line has one operator; \texttt{b * c} isn't allowed to appear as a subexpression
of the next line the way it could in the AST, it has to be materialized into
\texttt{t1} first.
\end{example}

\begin{example}[Translating an \texttt{if}]
\begin{lstlisting}
if (x > 10) { y = 1; } else { y = 2; }
\end{lstlisting}
becomes explicit comparisons, conditional jumps, and labels:
\begin{lstlisting}
t1 = x > 10
if t1 goto L_then
goto L_else
L_then:
    y = 1
    goto L_end
L_else:
    y = 2
L_end:
\end{lstlisting}
Notice the \texttt{if}/\texttt{else} \emph{structure} from the AST has disappeared,
replaced by explicit labels and gotos; TAC has no nested control-flow constructs at
all, just a flat sequence of instructions and jumps, which is exactly the shape most
optimizations want to work with.
\end{example}

\begin{example}[Translating a \texttt{while}]
\begin{lstlisting}
while (i < n) { s = s + i; i = i + 1; }
\end{lstlisting}
becomes:
\begin{lstlisting}
L_test:
    t1 = i < n
    if t1 goto L_body
    goto L_end
L_body:
    s = s + i
    i = i + 1
    goto L_test
L_end:
\end{lstlisting}
Same idea: the loop's structure is entirely encoded by labels and jumps, there's no
special ``loop'' instruction at all.
\end{example}

\subsection{SSA form}

\textbf{Static single assignment} (SSA) is a further refinement of an IR like TAC:
every variable is assigned to \emph{exactly once}, textually. If the source program
reassigns a variable, SSA gives each reassignment a fresh, versioned name instead.

The wrinkle is control flow merging: if a variable is assigned differently along two
branches, and then used after the branches merge back together, which assignment is
``the'' value at the merge point? SSA's answer is a special \textbf{phi function}
($\phi$), placed at the merge point, that says ``take the value from whichever branch
we actually came from''.

\begin{example}[Before/after SSA across a branch]
Before (ordinary TAC, \texttt{x} reassigned along both branches):
\begin{lstlisting}
if (cond) goto L1
x = 1
goto L2
L1:
x = 2
L2:
y = x + 1
\end{lstlisting}
After SSA (every assignment gets a fresh version, and the merge point gets a $\phi$):
\begin{lstlisting}
if (cond) goto L1
x1 = 1
goto L2
L1:
x2 = 2
L2:
x3 = phi(x1, x2)   ; x1 if we arrived from the fall-through, x2 if from L1
y1 = x3 + 1
\end{lstlisting}
Now every use of a variable (\texttt{x1}, \texttt{x2}, \texttt{x3}, \texttt{y1}) refers,
unambiguously, to exactly one assignment. That property is what makes SSA so friendly
to data-flow analyses: ``where did this value come from'' is always a direct,
syntactic answer, instead of a question that requires tracing execution paths
backward through possible reassignments.
\end{example}

\subsection{from IR to native code}

Getting from an IR like TAC (or its SSA-form refinement) down to actual machine
instructions has two big conceptual pieces.

\subsubsection{instruction selection}

\textbf{Instruction selection} is the problem of matching IR patterns to the target
machine's actual instructions. A single TAC instruction might correspond to one
machine instruction, several, or (if the target has a fused instruction for a common
pattern, a multiply-add, say) fewer than the naive one-to-one count would suggest.
This is a pattern-matching problem: find, for each chunk of IR, the cheapest sequence
of real instructions that computes the same thing.

\subsubsection{register allocation}

The IR happily invents as many temporaries as it wants (\texttt{t1}, \texttt{t2}, \ldots,
unboundedly many), but real hardware has a small, fixed number of registers (a few
dozen, at most, on most architectures). \textbf{Register allocation} is the problem of
mapping that unbounded set of IR temporaries onto the small fixed set of physical
registers, deciding, for each point in the program, which temporaries get to live in a
real register and which have to be \textbf{spilled} to memory (stored to and reloaded
from the stack) because there simply aren't enough registers to go around at that
point.

The classic conceptual framework here is \textbf{graph-coloring register allocation}:
build an \emph{interference graph} where each node is a temporary and an edge connects
two temporaries that are both live (still needed) at the same program point, so they
can't share a register. Assigning registers is then, quite literally, the graph-coloring
problem: color the graph with $k$ colors ($k$ = number of physical registers) such that
no two adjacent nodes share a color; a temporary that can't be colored has to spill. We
won't derive the full algorithm here, but it's worth knowing the name, because it's the
textbook answer whenever someone asks ``how does a real compiler decide what goes in a
register''.

\begin{remark}
Instruction selection and register allocation are, in principle, separate concerns
(matching \emph{what} instruction to use vs.\ deciding \emph{where}, register or
memory, a value lives), but real compilers often have to interleave decisions between
the two, since the best instruction choice can depend on register pressure and vice
versa. This is one of the genuinely hard, still-actively-researched parts of compiler
backend design.
\end{remark}

\subsection{Exercises}
\begin{itemize}
  \item Translate \texttt{z = (a + b) * (c - d)} into three-address code by hand,
    naming temporaries \texttt{t1}, \texttt{t2}, \ldots\ in the order they're
    introduced.
  \item Translate \texttt{if (a > b) \{ m = a; \} else \{ m = b; \}} into three-address
    code with explicit labels and jumps, following the pattern shown for \texttt{if}
    above.
  \item Convert your answer to the previous exercise into SSA form, introducing a
    $\phi$ function at the merge point. Which variable versions does the $\phi$
    combine, and under what condition does each one apply?
  \item Trace a small while-loop of your choice (2--3 statements in the body) all the
    way from source, to TAC, to SSA form, following the patterns shown in this
    section.
  \item Think about it: why does SSA form make an optimization like ``is this value
    ever used again after this point'' (dead code elimination) mechanically easier to
    check than it would be on ordinary, non-SSA TAC where a variable might be
    reassigned several times?
  \item Suppose a target machine has only 3 physical registers, and at some point in
    the program 4 different temporaries are simultaneously live (needed later).
    Describe, in your own words, what the register allocator has to do here, and what
    the runtime cost of that decision is compared to a hypothetical machine with
    enough registers for everything.
\end{itemize}
